\Needspace{10\baselineskip}\subsection*{CHEM: Chemical revolution}
\addcontentsline{toc}{subsection}{CHEM: Chemical revolution}
\noindent\textit{Predecessor} $T_1$: Phlogistic chemistry (Stahl; Black, Priestley, Cavendish, Scheele, Kirwan), c.~1700--1787.\\
\textit{Successor} $T_2$: Oxygen chemistry (Lavoisier, Trait\'{e} \'{e}l\'{e}mentaire de chimie, 1789) and its continuation.\par\smallskip
\begin{manualtable}
\textbf{CH01} & \textbf{phlogiston}. Principle of inflammability, contained in all combustibles and metals and released into the air in combustion and calcination. & --- & \textbf{T} token-dependent; $w$\,=\,10; $s$\,=\,1/4 & by token \\
\quad CH01a & Stahl and his followers: the principle common to all combustibles, given off when they burn & \textit{no referent} & fails & U: I01 I04 I09 I12 I20 I26 \\
\quad CH01b & In reduction: what a calx takes from charcoal when heated with it to restore the metal & \textit{no referent} & fails & U: I08 \\
\quad CH01c & Kirwan (Essay on Phlogiston, 1787) and some statements of Cavendish: phlogiston identified with, or as the main constituent of, the inflammable air obtained from metals and acids & hydrogen & refers & H: I28 \\
\quad CH01d & What air absorbs from burning bodies until it is saturated and combustion stops & \textit{no referent} & fails & U: I02 I22 \\
\midrule
\textbf{CH02} & \textbf{dephlogisticated air}. Air deprived of its phlogiston and therefore unusually avid of it; supports combustion and respiration better than common air. & --- & \textbf{T} token-dependent; $w$\,=\,4; $s$\,=\,2/3 & by token \\
\quad CH02a & Priestley, 1774--1775: the air obtained by heating the red calx of mercury, in which a candle burned with a remarkably vigorous flame & oxygen & refers & H: I10 I11 \\
\quad CH02b & The type description: common air from which phlogiston has been removed & \textit{no referent} & fails & U: I12 \\
\quad CH02c & Cavendish, 1784: the air exploded with inflammable air in his synthesis of water & oxygen & refers & H: I18 \\
\midrule
\textbf{CH03} & \textbf{phlogisticated air}. Air saturated with phlogiston, which no longer supports combustion or respiration. & --- & \textbf{T} token-dependent; $w$\,=\,2; $s$\,=\,1/2 & by token \\
\quad CH03a & Residue of air in which charcoal or a candle has burned, after the fixed air has been absorbed (Rutherford, 1772; Priestley) & nitrogen (with traces of the noble gases, then unknown) & refers & H: I14 \\
\quad CH03b & The type description: common air that has taken up as much phlogiston as it can hold & \textit{no referent} & fails & U: I02 \\
\midrule
\textbf{CH04} & \textbf{inflammable air}. An air that burns; for Cavendish (1766) the air obtained from metals dissolved in acids. & --- & \textbf{T} token-dependent; $w$\,=\,5; $s$\,=\,1 & by token \\
\quad CH04a & Cavendish, 1766: air from iron, zinc or tin dissolved in dilute vitriolic or marine acid & hydrogen & refers & H: I17 I18 I28 \\
\quad CH04b & Volta, 1776: inflammable air collected from the mud of marshes & methane & refers & H: I29 \\
\quad CH04c & The heavy inflammable air from heating charcoal with metallic calxes (Priestley) & carbon monoxide & refers & H: I30 \\
\midrule
\textbf{CH05} & \textbf{fixed air}. Air fixed in chalk and magnesia alba and released on heating or by acids; extinguishes flame. & carbon dioxide (carbonic acid gas) & \textbf{E} exact; $w$\,=\,3; $s$\,=\,1 & H: I14 I15 I16 \\
\quad CH05a & Black, 1756: the air expelled from magnesia alba and chalk by heat & carbon dioxide & refers & as term \\
\quad CH05b & Air given off in fermentation and exhaled in respiration & carbon dioxide & refers & as term \\
\midrule
\textbf{CH06} & \textbf{calx}. The earthy residue of the calcination of a metal; the metal deprived of its phlogiston. & metallic oxide & \textbf{P} partial; $w$\,=\,6; $s$\,=\,1 & H: I06 I07 I10 I30; U: I04 I08 \\
\quad CH06a & Red calx of mercury (mercurius calcinatus per se) & mercuric oxide & refers & as term \\
\quad CH06b & Calx of lead (litharge, minium) & lead oxides & refers & as term \\
\quad CH06c & Calx of tin formed in a sealed retort (Lavoisier, 1774) & tin oxide & refers & as term \\
\midrule
\textbf{CH07} & \textbf{metal}. A compound of a calx with phlogiston; the common properties of metals are due to their common phlogiston. & metal (a simple substance) & \textbf{P} partial; $w$\,=\,7; $s$\,=\,1 & H: I05 I06 I07 I10 I17; U: I04 I26 \\
\quad CH07a & Mercury, lead, tin, iron, zinc as named in calcination experiments & the same metals & refers & as term \\
\quad CH07b & Metals as a class sharing lustre and malleability & the metals & refers & as term \\
\midrule
\textbf{CH08} & \textbf{combustion}. Release of phlogiston from a combustible into the surrounding air. & combustion (combination with oxygen) & \textbf{P} partial; $w$\,=\,10; $s$\,=\,1 & H: I03 I11 I13 I14 I16 I23 I31; U: I01 I02 I09 \\
\quad CH08a & Burning of a candle, charcoal or sulphur in a closed volume of air & combustion in air & refers & as term \\
\midrule
\textbf{CH09} & \textbf{calcination}. Conversion of a metal into its calx by the loss of phlogiston on heating in air. & oxidation of a metal & \textbf{E} exact; $w$\,=\,2; $s$\,=\,1 & H: I05 I06 \\
\quad CH09a & Heating of tin or lead in air until it becomes a calx & oxidation of tin or lead & refers & as term \\
\midrule
\textbf{CH10} & \textbf{reduction (revivification)}. Restoration of a metal from its calx by restoring phlogiston, usually from charcoal. & reduction (removal of oxygen) & \textbf{P} partial; $w$\,=\,2; $s$\,=\,1 & H: I07; U: I08 \\
\quad CH10a & Heating a calx with charcoal to recover the metal & reduction of the oxide by carbon & refers & as term \\
\midrule
\textbf{CH11} & \textbf{common air}. The atmosphere, treated as one kind of air whose goodness varies with the phlogiston it holds. & atmospheric air (a mixture chiefly of oxygen and nitrogen) & \textbf{P} partial; $w$\,=\,6; $s$\,=\,1 & H: I03 I05 I11 I13; X: I27; U: I01 \\
\quad CH11a & The air of a room or of the open atmosphere used in eudiometry & atmospheric air & refers & as term \\
\midrule
\textbf{CH12} & \textbf{goodness of air}. Fitness of an air for respiration, inversely related to its content of phlogiston; measured by the nitrous air test. & proportion of oxygen in an air & \textbf{P} partial; $w$\,=\,2; $s$\,=\,1 & H: I21; U: I22 \\
\quad CH12a & Priestley's nitrous air test: the diminution of volume when nitrous air is mixed with the sample over water & oxygen content of the sample & refers & as term \\
\midrule
\textbf{CH13} & \textbf{nitrous air}. Air obtained from metals and spirit of nitre; diminishes good air when mixed with it over water. & nitric oxide & \textbf{E} exact; $w$\,=\,1; $s$\,=\,1 & H: I21 \\
\quad CH13a & Priestley, 1772: air from copper or brass dissolved in spirit of nitre & nitric oxide & refers & as term \\
\midrule
\textbf{CH14} & \textbf{dephlogisticated nitrous air}. Air obtained by exposing nitrous air to iron and other substances; supports a candle flame. & nitrous oxide & \textbf{E} exact; $w$\,=\,1; $s$\,=\,1 & H: I23 \\
\quad CH14a & Priestley's preparation; later the subject of Davy's researches (1800) & nitrous oxide & refers & as term \\
\midrule
\textbf{CH15} & \textbf{marine acid air}. Air obtained from spirit of salt; very soluble in water. & hydrogen chloride & \textbf{E} exact; $w$\,=\,1; $s$\,=\,1 & H: I24 \\
\quad CH15a & Priestley: air collected over mercury from heated spirit of salt & hydrogen chloride & refers & as term \\
\midrule
\textbf{CH16} & \textbf{dephlogisticated marine acid}. Marine acid deprived of phlogiston by manganese (Scheele, 1774); bleaches vegetable colours. & oxygenated muriatic acid (Lavoisierian name; now chlorine) & \textbf{E} exact; $w$\,=\,1; $s$\,=\,1 & H: I25 \\
\quad CH16a & Scheele: the yellow-green air from marine acid and manganese & chlorine & refers & as term \\
\midrule
\textbf{CH17} & \textbf{alkaline air}. Air obtained from sal ammoniac and lime; very soluble in water, alkaline. & ammoniacal gas & \textbf{E} exact; $w$\,=\,1; $s$\,=\,1 & H: I24 \\
\quad CH17a & Priestley: air collected over mercury from sal ammoniac and slaked lime & ammonia & refers & as term \\
\midrule
\textbf{CH18} & \textbf{vitriolic acid}. Oil of vitriol; in Stahl's chemistry, the acid which with phlogiston composes sulphur. & sulphuric acid & \textbf{P} partial; $w$\,=\,2; $s$\,=\,1 & H: I17; U: I20 \\
\quad CH18a & Oil of vitriol used to dissolve metals & sulphuric acid & refers & as term \\
\midrule
\textbf{CH19} & \textbf{sulphur}. A compound of vitriolic acid and phlogiston (Stahl). & sulphur (a simple substance) & \textbf{P} partial; $w$\,=\,2; $s$\,=\,1 & H: I31; U: I20 \\
\quad CH19a & Common brimstone burned in air & sulphur & refers & as term \\
\midrule
\textbf{CH20} & \textbf{charcoal}. A substance very rich in phlogiston, burning almost without residue. & carbon (as charcoal) & \textbf{P} partial; $w$\,=\,5; $s$\,=\,1 & H: I07 I16 I30; U: I08 I09 \\
\quad CH20a & Charcoal used in reduction and in combustion experiments & impure carbon & refers & as term \\
\midrule
\textbf{CH21} & \textbf{water}. An element; in Cavendish's 1784 interpretation, the base of both dephlogisticated and inflammable air. & water (a compound of hydrogen and oxygen) & \textbf{P} partial; $w$\,=\,3; $s$\,=\,1 & H: I18 I31; X: I19 \\
\quad CH21a & The dew deposited in Cavendish's exploding vessel & water & refers & as term \\
\midrule
\textbf{CH22} & \textbf{element}. One of the few ultimate constituents (earth, water, air, fire, or Stahl's principles), not identified by analysis. & simple substance (Lavoisier's operational sense: not yet decomposed) & \textbf{P} partial; $w$\,=\,2; $s$\,=\,1 & X: I19 I27 \\
\quad CH22a & Water and air counted among the elements & simple substance (the predicate) & refers & as term \\
\midrule
\textbf{CH23} & \textbf{respiration}. Discharge of phlogiston from the lungs into the air. & respiration (slow combustion consuming oxygen) & \textbf{E} exact; $w$\,=\,3; $s$\,=\,1 & H: I11 I13 I16 \\
\quad CH23a & Mice and men breathing confined air until it is unfit & respiration & refers & as term \\
\bottomrule
\end{manualtable}

\Needspace{10\baselineskip}\subsection*{ASTR: Astronomical revolution}
\addcontentsline{toc}{subsection}{ASTR: Astronomical revolution}
\noindent\textit{Predecessor} $T_1$: Ptolemaic astronomy (Almagest, c.~150; Planetary Hypotheses) with Aristotelian cosmology, as taught to 1543.\\
\textit{Successor} $T_2$: Copernican--Keplerian astronomy (De revolutionibus, 1543; Astronomia nova, 1609; Rudolphine Tables, 1627).\par\smallskip
\begin{manualtable}
\textbf{AS01} & \textbf{planet}. A wandering star: Moon, Mercury, Venus, Sun, Mars, Jupiter, Saturn. & --- & \textbf{T} token-dependent; $w$\,=\,13; $s$\,=\,1 & by token \\
\quad AS01a & Mars, Jupiter or Saturn called a planet & the same body, a planet & refers & H: I02 I03 I09 I21; R: I01 I05; U: I11 \\
\quad AS01b & The Sun called a planet & the Sun, a star at the centre of the system & refers & X: I20 I34 \\
\quad AS01c & The Moon called a planet & the Moon, a satellite of the Earth & refers & H: I26 \\
\quad AS01d & Venus or Mercury called a planet & the same body, a planet & refers & H: I19 I22; R: I04 \\
\midrule
\textbf{AS02} & \textbf{fixed star}. A star keeping its position relative to the other stars; carried round daily by the heavens. & star (beyond the planetary system) & \textbf{P} partial; $w$\,=\,3; $s$\,=\,1 & H: I09 I13; X: I14 \\
\quad AS02a & The stars catalogued in Almagest VII--VIII & the same stars & refers & as term \\
\midrule
\textbf{AS03} & \textbf{primum mobile}. The outermost moving sphere, which turns daily from east to west and carries the inner spheres with it. & --- & \textbf{N} none; $w$\,=\,1; $s$\,=\,0 & by token \\
\quad AS03a & The sphere that communicates the daily motion to all the heavens & \textit{no referent} & fails & U: I08 \\
\midrule
\textbf{AS04} & \textbf{planetary orb}. A solid spherical shell of aether carrying a planet; the orbs nest without gaps (Planetary Hypotheses; Peurbach's total and partial orbs). & --- & \textbf{N} none; $w$\,=\,4; $s$\,=\,0 & by token \\
\quad AS04a & The orb of Mars as a solid shell & \textit{no referent} & fails & U: I08 I11 I33 \\
\quad AS04b & The orbs whose nesting fixes absolute distances in the Planetary Hypotheses & \textit{no referent} & fails & U: I29 \\
\midrule
\textbf{AS05} & \textbf{deferent}. The large circle about the Earth (or an eccentric point) on which the centre of a planet's epicycle moves. & --- & \textbf{T} token-dependent; $w$\,=\,2; $s$\,=\,1 & by token \\
\quad AS05a & Deferent of a superior planet (Mars, Jupiter, Saturn) & the planet's orbit about the Sun & refers & R: I01 \\
\quad AS05b & Deferent of an inferior planet (Venus, Mercury), whose epicycle centre stays on the line to the mean Sun & the Earth's orbit about the Sun (seen from the Earth as the Sun's annual motion) & refers & R: I04 \\
\midrule
\textbf{AS06} & \textbf{epicycle}. A circle whose centre moves on the deferent and on which the planet moves; accounts for the second anomaly. & --- & \textbf{T} token-dependent; $w$\,=\,5; $s$\,=\,1 & by token \\
\quad AS06a & Main epicycle of a superior planet & the Earth's orbit about the Sun, as reflected in the planet's geocentric motion & refers & R: I01 I05 I15 \\
\quad AS06b & Main epicycle of Venus or Mercury & the planet's orbit about the Sun & refers & R: I04 \\
\quad AS06c & Epicycle of the Moon (first lunar model) & the eccentricity of the Moon's orbit about the Earth (the equation of centre) & refers & H: I28 \\
\midrule
\textbf{AS07} & \textbf{eccentric}. A circle whose centre is displaced from the Earth. & eccentricity of an orbit (displacement of the Sun from its centre) & \textbf{P} partial; $w$\,=\,2; $s$\,=\,1 & R: I16 I25 \\
\quad AS07a & Eccentric of the Sun's circle (Hipparchus, Almagest III) & the eccentricity of the Earth's orbit & refers & as term \\
\midrule
\textbf{AS08} & \textbf{equant}. The point about which the centre of the epicycle moves with uniform angular speed. & --- & \textbf{T} token-dependent; $w$\,=\,3; $s$\,=\,1/2 & by token \\
\quad AS08a & Equant of a superior planet, at twice the eccentricity of the deferent's centre & the empty focus of the planet's ellipse (to first order in the eccentricity) & refers & R: I15 I16 \\
\quad AS08b & The equant as the centre of a physically uniform rotation of the orb & \textit{no referent} & fails & U: I33 \\
\midrule
\textbf{AS09} & \textbf{first anomaly}. The non-uniformity of a planet's motion that depends on its position in the zodiac. & non-uniform orbital motion due to eccentricity (Kepler's second law) & \textbf{E} exact; $w$\,=\,1; $s$\,=\,1 & H: I18 \\
\quad AS09a & The zodiacal anomaly of Mars & the non-uniform heliocentric motion of Mars & refers & as term \\
\midrule
\textbf{AS10} & \textbf{second anomaly}. The non-uniformity, including retrogradation, that depends on the planet's position relative to the Sun. & the apparent effect of the Earth's orbital motion & \textbf{E} exact; $w$\,=\,2; $s$\,=\,1 & H: I17 I31 \\
\quad AS10a & The synodic anomaly of Jupiter & the parallax of the Earth's orbit in Jupiter's direction & refers & as term \\
\midrule
\textbf{AS11} & \textbf{retrograde motion}. A real backward motion of the planet, produced by its motion on the epicycle. & apparent retrograde motion & \textbf{P} partial; $w$\,=\,3; $s$\,=\,1 & H: I02 I03; X: I30 \\
\quad AS11a & Retrogradation of Mars near opposition & apparent retrograde motion of Mars & refers & as term \\
\midrule
\textbf{AS12} & \textbf{Earth}. The stationary body at the centre of the universe. & the Earth (a planet) & \textbf{P} partial; $w$\,=\,4; $s$\,=\,1 & R: I04 I05 I25; X: I06 \\
\quad AS12a & The Earth on which observers stand & the Earth & refers & as term \\
\midrule
\textbf{AS13} & \textbf{Sun}. A planet moving about the Earth on an eccentric circle. & the Sun & \textbf{P} partial; $w$\,=\,9; $s$\,=\,1 & H: I02 I22 I24 I27; R: I04 I05 I25; X: I20 I34 \\
\quad AS13a & The Sun whose annual motion defines the ecliptic & the Sun & refers & as term \\
\midrule
\textbf{AS14} & \textbf{Moon}. The planet nearest the Earth. & the Moon (the Earth's satellite) & \textbf{P} partial; $w$\,=\,5; $s$\,=\,1 & H: I26 I28 I32; X: I20 I23 \\
\quad AS14a & The Moon of eclipses and lunar theory & the Moon & refers & as term \\
\midrule
\textbf{AS15} & \textbf{centre of the universe}. The place to which heavy bodies tend and about which the heavens rotate; the centre of the Earth. & --- & \textbf{T} token-dependent; $w$\,=\,2; $s$\,=\,1/2 & by token \\
\quad AS15a & The place to which heavy bodies naturally move & the centre of the Earth & refers & H: I07 \\
\quad AS15b & The fixed point about which the heavens make their daily rotation & \textit{no referent} & fails & U: I10 \\
\midrule
\textbf{AS16} & \textbf{diurnal motion}. The daily rotation of the whole heavens from east to west. & apparent diurnal motion (due to the Earth's rotation) & \textbf{P} partial; $w$\,=\,3; $s$\,=\,1 & H: I09; U: I08 I10 \\
\quad AS16a & The daily rising and setting of all celestial bodies & apparent diurnal motion & refers & as term \\
\midrule
\textbf{AS17} & \textbf{ecliptic}. The great circle traced annually by the Sun among the fixed stars. & ecliptic (projection of the plane of the Earth's orbit) & \textbf{E} exact; $w$\,=\,2; $s$\,=\,1 & H: I27 I32 \\
\quad AS17a & The circle of the zodiac in which eclipses occur & ecliptic & refers & as term \\
\midrule
\textbf{AS18} & \textbf{tropical year}. The period of the Sun's return to the same equinox. & tropical year & \textbf{E} exact; $w$\,=\,1; $s$\,=\,1 & H: I24 \\
\quad AS18a & Ptolemy's value, slightly short of 365 1/4 days & tropical year & refers & as term \\
\midrule
\textbf{AS19} & \textbf{precession of the equinoxes}. A slow eastward motion of the sphere of fixed stars relative to the equinoxes. & precession (slow motion of the Earth's axis) & \textbf{P} partial; $w$\,=\,2; $s$\,=\,1 & H: I24; X: I14 \\
\quad AS19a & Hipparchus' and Ptolemy's precession & precession & refers & as term \\
\midrule
\textbf{AS20} & \textbf{synodic period}. The period of the second anomaly: return to the same configuration with the Sun. & synodic period & \textbf{E} exact; $w$\,=\,1; $s$\,=\,1 & H: I17 \\
\quad AS20a & Synodic period of Mars (about 780 days) & synodic period of Mars & refers & as term \\
\midrule
\textbf{AS21} & \textbf{zodiacal period}. The period of the mean motion of the epicycle centre around the zodiac. & --- & \textbf{T} token-dependent; $w$\,=\,2; $s$\,=\,1 & by token \\
\quad AS21a & Zodiacal period of a superior planet (Saturn about 29 1/2 years) & sidereal period of the planet & refers & H: I21 \\
\quad AS21b & Zodiacal period of Venus or Mercury (one year) & sidereal period of the Earth & refers & H: I19 \\
\midrule
\textbf{AS22} & \textbf{order of the planets}. The order by distance from the Earth, conventionally fixed by period: Moon, Mercury, Venus, Sun, Mars, Jupiter, Saturn. & order of distances (from the Earth or the Sun) & \textbf{P} partial; $w$\,=\,1; $s$\,=\,1 & X: I20 \\
\quad AS22a & The order stated in Almagest IX.1 & order of distances & refers & as term \\
\midrule
\textbf{AS23} & \textbf{elongation}. Angular distance of a planet from the Sun. & elongation & \textbf{E} exact; $w$\,=\,2; $s$\,=\,1 & H: I22 I31 \\
\quad AS23a & Greatest elongations of Venus and Mercury & elongation & refers & as term \\
\midrule
\textbf{AS24} & \textbf{aether}. The fifth element, of which the heavens are made: ungenerated, unalterable, moving in circles. & --- & \textbf{N} none; $w$\,=\,1; $s$\,=\,0 & by token \\
\quad AS24a & The substance of the orbs and stars & \textit{no referent} & fails & U: I12 \\
\midrule
\textbf{AS25} & \textbf{comet}. A burning exhalation in the upper sublunary region. & comet (a body beyond the Moon) & \textbf{P} partial; $w$\,=\,1; $s$\,=\,1 & X: I23 \\
\quad AS25a & The comet of 1577 & the comet of 1577 & refers & as term \\
\bottomrule
\end{manualtable}

\Needspace{10\baselineskip}\subsection*{MECH: Relativistic revolution}
\addcontentsline{toc}{subsection}{MECH: Relativistic revolution}
\noindent\textit{Predecessor} $T_1$: Newtonian mechanics and gravitation (Principia, 1687), with the ether-based electrodynamics of the late nineteenth century.\\
\textit{Successor} $T_2$: Special relativity (Einstein, 1905) and general relativity (Einstein, 1915).\par\smallskip
\begin{manualtable}
\textbf{ME01} & \textbf{mass}. Quantity of matter; invariant, additive and conserved; measures inertia and gravitational charge. & rest mass $\mid$ relativistic mass (total energy divided by $c^2$) (pd) & \textbf{P} partial; $w$\,=\,7; $s$\,=\,1 & H: I22; R: I01 I02 I03 I04 I26; U: I07 \\
\quad ME01a & Mass of a planet in the Kepler problem & rest mass $\approx$ relativistic mass (slow motion) & refers & as term \\
\quad ME01b & Mass of a fast electron in the deflection experiments of Kaufmann (1901--1906) & relativistic mass, or rest mass, as the reading is chosen & refers & as term \\
\quad ME01c & Mass of a composite body as the sum of the masses of its parts & relativistic mass (total energy over $c^2$) & refers & as term \\
\midrule
\textbf{ME02} & \textbf{force (impressed)}. Action exerted on a body changing its state of motion; equal to the rate of change of momentum. & --- & \textbf{T} token-dependent; $w$\,=\,3; $s$\,=\,2/3 & by token \\
\quad ME02a & Electric force on a charged particle & the Lorentz force (rate of change of relativistic momentum) & refers & H: I06 \\
\quad ME02b & Gravitational force of the Sun on a planet & \textit{no referent} & fails & U: I07 I28 \\
\quad ME02c & Contact force of a support or a spring & the same force (four-force in the rest frame) & refers & (none listed) \\
\midrule
\textbf{ME03} & \textbf{momentum (quantity of motion)}. Product of mass and velocity; conserved in isolated systems. & $m_0 v$ (low velocities) $\mid$ relativistic momentum $\gamma m_0 v$ (pd) & \textbf{P} partial; $w$\,=\,3; $s$\,=\,1 & H: I05 I06; R: I04 \\
\quad ME03a & Momentum of colliding billiard balls & $m_0 v \approx \gamma m_0 v$ & refers & as term \\
\quad ME03b & Momentum of a cathode-ray electron & relativistic momentum & refers & as term \\
\midrule
\textbf{ME04} & \textbf{kinetic energy (vis viva)}. One half of mass times the square of velocity. & kinetic energy $(\gamma-1)m_0c^2$ & \textbf{P} partial; $w$\,=\,1; $s$\,=\,1 & R: I26 \\
\quad ME04a & Energy of a falling body & kinetic energy & refers & as term \\
\midrule
\textbf{ME05} & \textbf{velocity}. Rate of change of position, relative to a frame; composes by vector addition. & velocity relative to an inertial frame & \textbf{P} partial; $w$\,=\,4; $s$\,=\,1 & R: I04 I10 I26; X: I19 \\
\quad ME05a & Velocity of a ship relative to the shore & velocity relative to a frame & refers & as term \\
\midrule
\textbf{ME06} & \textbf{absolute space}. Space in its own nature, without relation to anything external, always similar and immovable. & --- & \textbf{T} token-dependent; $w$\,=\,3; $s$\,=\,1/3 & by token \\
\quad ME06a & The immovable space of the Scholium to the Definitions & \textit{no referent} & fails & U: I13 \\
\quad ME06b & That relative to which the water in the rotating bucket rotates, as shown by its concave surface & the local inertial (affine) structure of spacetime & refers & H: I14 \\
\quad ME06c & The rest frame of the centre of the world (Principia, Book III, Hypothesis I) & \textit{no referent} & fails & U: I35 \\
\midrule
\textbf{ME07} & \textbf{absolute time}. Time that flows equably without relation to anything external. & --- & \textbf{T} token-dependent; $w$\,=\,2; $s$\,=\,1/2 & by token \\
\quad ME07a & True and mathematical time of the Scholium & \textit{no referent} & fails & U: I15 \\
\quad ME07b & The time variable in the equations of planetary motion & proper time along the observer's worldline (to the accuracy of the Newtonian limit) & refers & R: I16 \\
\midrule
\textbf{ME08} & \textbf{simultaneity}. A relation between events holding or failing absolutely. & simultaneity relative to an inertial frame & \textbf{P} partial; $w$\,=\,2; $s$\,=\,1 & R: I09 I16 \\
\quad ME08a & Two distant lightning strikes judged simultaneous & simultaneity in the observer's rest frame & refers & as term \\
\midrule
\textbf{ME09} & \textbf{inertial frame}. A frame at rest or in uniform motion relative to absolute space. & global inertial frame of special relativity $\mid$ local freely falling frame of general relativity (pd) & \textbf{P} partial; $w$\,=\,4; $s$\,=\,1 & H: I12 I34; R: I11 I30 \\
\quad ME09a & A frame at rest relative to the fixed stars, gravitation ignored & global inertial frame & refers & as term \\
\quad ME09b & A system of bodies falling freely in a uniform field (Principia, Corollary VI) & local freely falling frame & refers & as term \\
\midrule
\textbf{ME10} & \textbf{Galilean transformation}. The transformation between inertial frames: $x' = x - vt$, $t' = t$. & Lorentz transformation (of which it is the limit $v/c \to 0$) & \textbf{P} partial; $w$\,=\,2; $s$\,=\,1 & R: I10 I11 \\
\quad ME10a & Transformation from shore to ship frame & Lorentz transformation & refers & as term \\
\midrule
\textbf{ME11} & \textbf{gravitation}. A universal attractive force between masses, varying inversely as the square of the distance. & spacetime curvature produced by energy and momentum & \textbf{P} partial; $w$\,=\,5; $s$\,=\,1 & H: I22 I31; R: I21; U: I07 I23 \\
\quad ME11a & The gravitation that keeps the Moon in its orbit & curvature of spacetime near the Earth & refers & as term \\
\midrule
\textbf{ME12} & \textbf{gravitational potential}. The scalar whose gradient gives the gravitational acceleration. & the deviation of $g_{00}$ from its flat value (weak field) & \textbf{P} partial; $w$\,=\,1; $s$\,=\,1 & R: I08 \\
\quad ME12a & Potential of the Sun at the orbit of a planet & $(g_{00}+1)c^2/2$ in the weak field & refers & as term \\
\midrule
\textbf{ME13} & \textbf{orbit}. The path of a body under gravitation; a fixed conic for two bodies. & spatial path of a timelike geodesic & \textbf{P} partial; $w$\,=\,1; $s$\,=\,1 & R: I21 \\
\quad ME13a & The orbit of Mercury & the path of Mercury & refers & as term \\
\midrule
\textbf{ME14} & \textbf{length}. The extension of a body, independent of its motion. & length in a frame (proper length in the rest frame) & \textbf{P} partial; $w$\,=\,2; $s$\,=\,1 & H: I25; R: I17 \\
\quad ME14a & Length of a measuring rod & length of the rod in a frame & refers & as term \\
\midrule
\textbf{ME15} & \textbf{speed of light}. The speed of light relative to the ether, about 300\,000 km/s. & the invariant speed $c$ & \textbf{P} partial; $w$\,=\,4; $s$\,=\,1 & H: I20; X: I19; U: I18 I33 \\
\quad ME15a & R\o mer's and Fizeau's measured speed & the invariant speed $c$ & refers & as term \\
\midrule
\textbf{ME16} & \textbf{ether}. The medium of electromagnetic waves, at rest in absolute space. & --- & \textbf{T} token-dependent; $w$\,=\,3; $s$\,=\,1/3 & by token \\
\quad ME16a & A mechanical medium with elastic properties & \textit{no referent} & fails & U: I18 \\
\quad ME16b & Lorentz: the seat of the electric and magnetic fields & the electromagnetic field & refers & H: I20 \\
\quad ME16c & The rest frame relative to which the Earth's motion was sought by Michelson and Morley (1887) & \textit{no referent} & fails & U: I33 \\
\midrule
\textbf{ME17} & \textbf{gravitational constant}. The universal constant $G$ fixing the strength of gravitation. & the constant $G$ of the field equations & \textbf{E} exact; $w$\,=\,1; $s$\,=\,1 & H: I31 \\
\quad ME17a & Cavendish's torsion-balance value (1798) & the same constant & refers & as term \\
\midrule
\textbf{ME18} & \textbf{centre of mass}. The mass-weighted mean position of a system. & centre of energy in a frame & \textbf{E} exact; $w$\,=\,1; $s$\,=\,1 & H: I29 \\
\quad ME18a & Centre of mass of the solar system & centre of energy of the solar system & refers & as term \\
\midrule
\textbf{ME19} & \textbf{weight}. The gravitational force on a body. & --- & \textbf{T} token-dependent; $w$\,=\,2; $s$\,=\,1/2 & by token \\
\quad ME19a & The weight of a stone in the hand & \textit{no referent} & fails & U: I28 \\
\quad ME19b & What a spring balance registers & the supporting force preventing geodesic motion & refers & H: I32 \\
\midrule
\textbf{ME20} & \textbf{energy}. The conserved quantity comprising kinetic and potential energy. & energy (time component of four-momentum) & \textbf{E} exact; $w$\,=\,1; $s$\,=\,1 & H: I27 \\
\quad ME20a & Energy of a pendulum & energy & refers & as term \\
\midrule
\textbf{ME21} & \textbf{acceleration}. Rate of change of velocity; the same in all inertial frames. & coordinate acceleration $\mid$ proper acceleration (pd) & \textbf{P} partial; $w$\,=\,3; $s$\,=\,1 & H: I34; R: I08 I30 \\
\quad ME21a & Acceleration of a falling body & coordinate acceleration in the Earth frame & refers & as term \\
\quad ME21b & Acceleration felt by a passenger & proper acceleration & refers & as term \\
\midrule
\textbf{ME22} & \textbf{rigid body}. A body whose parts keep their distances; transmits impulses instantaneously. & --- & \textbf{T} token-dependent; $w$\,=\,2; $s$\,=\,1/2 & by token \\
\quad ME22a & Perfect rigidity as an idealization of transmission of impulse & \textit{no referent} & fails & U: I24 \\
\quad ME22b & A measuring rod moved slowly & a body in Born-rigid motion & refers & H: I25 \\
\midrule
\textbf{ME23} & \textbf{action at a distance}. Instantaneous action of one body on another without intervening medium. & --- & \textbf{N} none; $w$\,=\,1; $s$\,=\,0 & by token \\
\quad ME23a & Gravitational attraction across the solar system & \textit{no referent} & fails & U: I23 \\
\bottomrule
\end{manualtable}

\Needspace{10\baselineskip}\subsection*{HEAT: Thermodynamic revolution}
\addcontentsline{toc}{subsection}{HEAT: Thermodynamic revolution}
\noindent\textit{Predecessor} $T_1$: The caloric theory (Black's calorimetry as read by Lavoisier and Laplace, Laplace on sound, Carnot 1824), c.~1780--1840.\\
\textit{Successor} $T_2$: Thermodynamics and the kinetic theory of gases (Joule, Clausius, Kelvin, Maxwell, Boltzmann), 1840--1880.\par\smallskip
\begin{manualtable}
\textbf{HE01} & \textbf{caloric}. A subtle, elastic, weightless fluid whose particles repel one another and are attracted by ordinary matter; conserved. & --- & \textbf{T} token-dependent; $w$\,=\,8; $s$\,=\,1/2 & by token \\
\quad HE01a & The fluid of heat as described by Lavoisier and Laplace & \textit{no referent} & fails & U: I01 I08 I15 I18 I21 \\
\quad HE01b & The quantity of caloric given up by a body in the ice calorimeter & heat transferred (energy) & refers & H: I27 \\
\quad HE01c & Caloric as conserved in all processes & \textit{no referent} & fails & U: I02 \\
\quad HE01d & Carnot: the caloric that falls from the hot to the cold body in a reversible engine & entropy transferred (Callendar's reading); heat (Kelvin's) & refers & R: I14 \\
\midrule
\textbf{HE02} & \textbf{heat (quantity of)}. Quantity of caloric transferred. & heat (energy transferred by thermal interaction) & \textbf{P} partial; $w$\,=\,7; $s$\,=\,1 & H: I03 I04 I07 I23 I24; X: I13 I28 \\
\quad HE02a & Heat given by hot water to cold in Black's method of mixtures & heat transferred & refers & as term \\
\quad HE02b & Heat taken from the furnace by a steam engine & heat transferred & refers & as term \\
\midrule
\textbf{HE03} & \textbf{temperature}. Degree of heat; the tension or intensity of free caloric. & temperature & \textbf{P} partial; $w$\,=\,11; $s$\,=\,1 & H: I03 I04 I05 I07 I09 I12 I16 I19 I30; R: I14 I20 \\
\quad HE03a & Reading of a mercury thermometer & empirical temperature & refers & as term \\
\quad HE03b & Temperature of the boiler and condenser in Carnot & temperature (thermodynamic) & refers & as term \\
\midrule
\textbf{HE04} & \textbf{latent heat}. Caloric absorbed by a body in melting or vaporization, combined with it so as not to affect the thermometer. & latent heat (energy of phase change) & \textbf{P} partial; $w$\,=\,4; $s$\,=\,1 & H: I05 I24 I25; U: I06 \\
\quad HE04a & Black's latent heat of fusion of ice & latent heat of fusion & refers & as term \\
\quad HE04b & Latent heat of steam & latent heat of vaporization & refers & as term \\
\midrule
\textbf{HE05} & \textbf{specific heat}. The capacity of a body for caloric, per unit weight. & specific heat capacity & \textbf{P} partial; $w$\,=\,4; $s$\,=\,1 & H: I07 I11; U: I08 I15 \\
\quad HE05a & Specific heats measured by mixtures (Black, Crawford) & specific heat capacity & refers & as term \\
\midrule
\textbf{HE06} & \textbf{free caloric}. Caloric not combined with matter, sensible to the thermometer. & --- & \textbf{N} none; $w$\,=\,1; $s$\,=\,0 & by token \\
\quad HE06a & The free caloric whose density the thermometer measures & \textit{no referent} & fails & U: I26 \\
\midrule
\textbf{HE07} & \textbf{combined caloric}. Caloric held by chemical affinity in a body, insensible to the thermometer. & --- & \textbf{N} none; $w$\,=\,2; $s$\,=\,0 & by token \\
\quad HE07a & Caloric combined in water to form steam & \textit{no referent} & fails & U: I06 I29 \\
\midrule
\textbf{HE08} & \textbf{conduction}. Passage of heat through a body from hotter to colder parts. & heat conduction & \textbf{E} exact; $w$\,=\,1; $s$\,=\,1 & H: I16 \\
\quad HE08a & Fourier's heat flow in a solid (1822) & heat conduction & refers & as term \\
\midrule
\textbf{HE09} & \textbf{radiant heat}. Caloric emitted in free flight from bodies. & thermal radiation & \textbf{P} partial; $w$\,=\,2; $s$\,=\,1 & H: I17; U: I18 \\
\quad HE09a & Prevost's exchanges of radiant heat (1791) & thermal radiation & refers & as term \\
\quad HE09b & Herschel's rays beyond the red (1800) & infrared radiation & refers & as term \\
\midrule
\textbf{HE10} & \textbf{adiabatic compression}. Compression without exchange of caloric with the surroundings, which squeezes caloric out of combination and raises temperature. & adiabatic compression & \textbf{E} exact; $w$\,=\,3; $s$\,=\,1 & H: I09 I10 I27 \\
\quad HE10a & Compression in a sound wave (Laplace) & adiabatic compression & refers & as term \\
\midrule
\textbf{HE11} & \textbf{ratio of specific heats}. The ratio of the specific heat of a gas at constant pressure to that at constant volume. & ratio of specific heats $\gamma$ & \textbf{E} exact; $w$\,=\,2; $s$\,=\,1 & H: I10 I11 \\
\quad HE11a & The ratio inferred by Laplace from the speed of sound & $\gamma$ & refers & as term \\
\midrule
\textbf{HE12} & \textbf{speed of sound}. Speed of a pressure wave in air. & speed of sound & \textbf{E} exact; $w$\,=\,1; $s$\,=\,1 & H: I10 \\
\quad HE12a & The measured speed of sound in air & speed of sound & refers & as term \\
\midrule
\textbf{HE13} & \textbf{heat engine}. A device that produces motive power from the passage of caloric from a hot to a cold body. & heat engine & \textbf{P} partial; $w$\,=\,3; $s$\,=\,1 & H: I12; R: I14; X: I13 \\
\quad HE13a & The steam engine & heat engine & refers & as term \\
\midrule
\textbf{HE14} & \textbf{reversible process}. A process that can be run backward through the same states (Carnot). & reversible process & \textbf{E} exact; $w$\,=\,1; $s$\,=\,1 & H: I12 \\
\quad HE14a & Carnot's ideal cycle & Carnot cycle & refers & as term \\
\midrule
\textbf{HE15} & \textbf{frictional heating}. Release of caloric from bodies whose capacity is reduced by friction or abrasion. & conversion of work into heat by friction & \textbf{P} partial; $w$\,=\,2; $s$\,=\,1 & X: I28; U: I15 \\
\quad HE15a & Rumford's boring of cannon (1798) & conversion of work into heat & refers & as term \\
\midrule
\textbf{HE16} & \textbf{thermal equilibrium}. State in which bodies in contact no longer exchange caloric. & thermal equilibrium & \textbf{E} exact; $w$\,=\,2; $s$\,=\,1 & H: I17 I19 \\
\quad HE16a & Bodies left in contact & thermal equilibrium & refers & as term \\
\midrule
\textbf{HE17} & \textbf{thermal expansion}. Increase of volume due to the repulsion of added caloric. & thermal expansion & \textbf{P} partial; $w$\,=\,1; $s$\,=\,1 & R: I20 \\
\quad HE17a & Expansion of gases measured by Gay-Lussac (1802) & thermal expansion & refers & as term \\
\midrule
\textbf{HE18} & \textbf{elastic fluid (gas)}. A fluid whose particles are held apart by atmospheres of caloric. & gas & \textbf{P} partial; $w$\,=\,5; $s$\,=\,1 & H: I09 I11; R: I20 I22; U: I21 \\
\quad HE18a & Air and the permanent gases & gases & refers & as term \\
\midrule
\textbf{HE19} & \textbf{elasticity (of a gas)}. The spring of a gas, due to the repulsion of caloric atmospheres. & pressure & \textbf{P} partial; $w$\,=\,2; $s$\,=\,1 & R: I22; U: I21 \\
\quad HE19a & The spring of air in Boyle's experiments & pressure & refers & as term \\
\midrule
\textbf{HE20} & \textbf{vapour}. A liquid rendered elastic by combination with caloric. & vapour (the gaseous phase) & \textbf{P} partial; $w$\,=\,2; $s$\,=\,1 & H: I25; U: I29 \\
\quad HE20a & Steam & water vapour & refers & as term \\
\midrule
\textbf{HE21} & \textbf{calorimeter}. An instrument measuring caloric released by a body. & calorimeter & \textbf{E} exact; $w$\,=\,1; $s$\,=\,1 & H: I24 \\
\quad HE21a & Lavoisier and Laplace's ice calorimeter (1783) & ice calorimeter & refers & as term \\
\midrule
\textbf{HE22} & \textbf{thermometer}. An instrument measuring the degree of heat. & thermometer & \textbf{P} partial; $w$\,=\,2; $s$\,=\,1 & H: I30; U: I26 \\
\quad HE22a & Mercury thermometer & thermometer & refers & as term \\
\midrule
\textbf{HE23} & \textbf{animal heat}. Caloric released in respiration. & body heat (produced by metabolism) & \textbf{E} exact; $w$\,=\,1; $s$\,=\,1 & H: I23 \\
\quad HE23a & Heat given off by a guinea pig in the ice calorimeter (Lavoisier and Laplace) & metabolic heat & refers & as term \\
\bottomrule
\end{manualtable}

\Needspace{10\baselineskip}\subsection*{QUAN: Quantum revolution}
\addcontentsline{toc}{subsection}{QUAN: Quantum revolution}
\noindent\textit{Predecessor} $T_1$: Classical mechanics and electrodynamics applied to atoms and radiation, with the nuclear atom (Rutherford, 1911), c.~1895--1912.\\
\textit{Successor} $T_2$: Quantum mechanics (Heisenberg 1925; Born, Jordan; Schr\"odinger 1926; Dirac) in its standard formulation.\par\smallskip
\begin{manualtable}
\textbf{QU01} & \textbf{trajectory}. The continuous path of a particle, fixed by its position and momentum at one time and the forces. & --- & \textbf{T} token-dependent; $w$\,=\,4; $s$\,=\,2/3 & by token \\
\quad QU01a & Trajectory of an electron inside an atom & \textit{no referent} & fails & U: I30 \\
\quad QU01b & Track of an alpha particle in a cloud chamber (C.T.R. Wilson, 1911--1912) & the path of a well-localized wave packet (Mott, 1929) & refers & H: I09 \\
\quad QU01c & Trajectory of a projectile or planet & the path of the centre of a localized state (Ehrenfest's theorem) & refers & R: I08 I29 \\
\midrule
\textbf{QU02} & \textbf{electron orbit}. The closed path of an electron revolving about the nucleus. & --- & \textbf{T} token-dependent; $w$\,=\,4; $s$\,=\,1/2 & by token \\
\quad QU02a & Orbit of the electron in the ground state of hydrogen & \textit{no referent} & fails & U: I01 I02 I04 \\
\quad QU02b & Orbit of an electron in a very highly excited atom & a Rydberg wave packet circulating with the classical orbital period & refers & R: I05 \\
\midrule
\textbf{QU03} & \textbf{position}. A quantity possessed with a definite value by every particle at every time. & position (an observable) & \textbf{P} partial; $w$\,=\,3; $s$\,=\,1 & R: I08; X: I06 I07 \\
\quad QU03a & Position of an electron on a scintillation screen & position observable (measured value) & refers & as term \\
\midrule
\textbf{QU04} & \textbf{momentum}. Mass times velocity, possessed with a definite value at every time. & momentum (an observable) & \textbf{P} partial; $w$\,=\,3; $s$\,=\,1 & R: I08; X: I06 I07 \\
\quad QU04a & Momentum of a cathode-ray electron deflected in a magnetic field & momentum observable & refers & as term \\
\midrule
\textbf{QU05} & \textbf{energy}. A quantity taking any value in a continuous range for a bound system. & energy (the Hamiltonian, with discrete eigenvalues for bound states) & \textbf{P} partial; $w$\,=\,3; $s$\,=\,1 & R: I11; X: I10 I16 \\
\quad QU05a & Energy of the electron in an atom & energy of the bound state & refers & as term \\
\midrule
\textbf{QU06} & \textbf{electron}. J.J. Thomson's corpuscle (1897): a small charged particle with definite charge and mass. & electron & \textbf{P} partial; $w$\,=\,4; $s$\,=\,1 & H: I09 I17; X: I10; U: I01 \\
\quad QU06a & Cathode-ray corpuscles & electrons & refers & as term \\
\quad QU06b & Electrons in Millikan's oil drops & electrons & refers & as term \\
\midrule
\textbf{QU07} & \textbf{atom}. A nucleus with electrons moving about it under electric forces (Rutherford). & atom & \textbf{P} partial; $w$\,=\,4; $s$\,=\,1 & H: I19 I28; X: I10; U: I01 \\
\quad QU07a & The hydrogen atom & hydrogen atom & refers & as term \\
\midrule
\textbf{QU08} & \textbf{nucleus}. A small, massive, positively charged centre of the atom. & nucleus & \textbf{P} partial; $w$\,=\,5; $s$\,=\,1 & H: I19 I20; U: I01 I02 I30 \\
\quad QU08a & Rutherford's nucleus (1911) & nucleus & refers & as term \\
\midrule
\textbf{QU09} & \textbf{light wave}. An electromagnetic wave in which energy is continuously distributed over the wave front. & electromagnetic radiation (the quantized field) & \textbf{P} partial; $w$\,=\,5; $s$\,=\,1 & H: I15 I21 I27; X: I14 I16 \\
\quad QU09a & Light in an interference experiment & electromagnetic radiation & refers & as term \\
\midrule
\textbf{QU10} & \textbf{radiation of an accelerated charge}. Continuous emission of electromagnetic energy by every accelerated charge (Larmor, 1897). & electromagnetic emission by charges & \textbf{P} partial; $w$\,=\,3; $s$\,=\,1 & H: I27; R: I03; U: I02 \\
\quad QU10a & Radiation from currents in an antenna & emission from macroscopic currents & refers & as term \\
\quad QU10b & Radiation expected from an electron revolving in an atom & transition radiation & refers & as term \\
\midrule
\textbf{QU11} & \textbf{spectral line}. Radiation emitted at the frequency of a periodic motion of the electrons or its harmonics. & transition frequency between stationary states & \textbf{P} partial; $w$\,=\,3; $s$\,=\,1 & H: I28; R: I05; U: I04 \\
\quad QU11a & The Balmer lines of hydrogen & transition frequencies of hydrogen & refers & as term \\
\midrule
\textbf{QU12} & \textbf{equipartition}. In thermal equilibrium each degree of freedom has mean energy $kT/2$. & the high-temperature limit of the quantum mean energy & \textbf{P} partial; $w$\,=\,3; $s$\,=\,1 & R: I11 I12 I13 \\
\quad QU12a & Equipartition among vibrational modes of a solid & high-temperature limit & refers & as term \\
\midrule
\textbf{QU13} & \textbf{black-body radiation}. Radiation in thermal equilibrium within an enclosure. & black-body radiation & \textbf{P} partial; $w$\,=\,1; $s$\,=\,1 & R: I12 \\
\quad QU13a & Cavity radiation measured by Lummer and Pringsheim, Rubens and Kurlbaum & black-body radiation & refers & as term \\
\midrule
\textbf{QU14} & \textbf{atomic heat}. The product of specific heat and atomic weight of a solid element, constant by the law of Dulong and Petit. & molar heat capacity of a solid & \textbf{P} partial; $w$\,=\,1; $s$\,=\,1 & R: I13 \\
\quad QU14a & Atomic heat of copper at room temperature & molar heat capacity & refers & as term \\
\midrule
\textbf{QU15} & \textbf{photoelectric emission}. Release of electrons from metals by light (Hertz, 1887; Lenard, 1902). & photoelectric effect & \textbf{P} partial; $w$\,=\,1; $s$\,=\,1 & X: I14 \\
\quad QU15a & Lenard's measurements of emitted electron energies & photoelectric effect & refers & as term \\
\midrule
\textbf{QU16} & \textbf{angular momentum}. Moment of momentum, taking any value. & angular momentum (quantized observable; includes spin) & \textbf{P} partial; $w$\,=\,3; $s$\,=\,1 & H: I23; X: I22 I24 \\
\quad QU16a & Angular momentum of an orbiting electron & orbital angular momentum & refers & as term \\
\midrule
\textbf{QU17} & \textbf{magnetic moment}. The moment of a current loop, taking any direction relative to a field. & magnetic moment (quantized components) & \textbf{P} partial; $w$\,=\,1; $s$\,=\,1 & X: I24 \\
\quad QU17a & Magnetic moment of a silver atom (Stern and Gerlach, 1922) & magnetic moment & refers & as term \\
\midrule
\textbf{QU18} & \textbf{state (mechanical)}. A point in phase space: positions and momenta of all particles. & --- & \textbf{T} token-dependent; $w$\,=\,2; $s$\,=\,1/2 & by token \\
\quad QU18a & The microstate of a system of electrons & \textit{no referent} & fails & U: I25 \\
\quad QU18b & The macroscopic state of a gas given by pressure, volume and temperature & the thermodynamic macrostate & refers & H: I26 \\
\midrule
\textbf{QU19} & \textbf{charge}. Electric charge, conserved and coming in multiples of $e$. & charge & \textbf{P} partial; $w$\,=\,3; $s$\,=\,1 & H: I17 I18; R: I03 \\
\quad QU19a & Millikan's elementary charge (1909--1913) & elementary charge & refers & as term \\
\midrule
\textbf{QU20} & \textbf{X-rays}. Short-wavelength electromagnetic radiation. & X-rays & \textbf{E} exact; $w$\,=\,1; $s$\,=\,1 & H: I21 \\
\quad QU20a & X-rays diffracted by crystals (Laue, 1912) & X-rays & refers & as term \\
\midrule
\textbf{QU21} & \textbf{Rutherford scattering}. Deflection of alpha particles by the Coulomb field of a nucleus. & Coulomb scattering & \textbf{E} exact; $w$\,=\,1; $s$\,=\,1 & H: I20 \\
\quad QU21a & Geiger and Marsden's scattering (1909--1913) & Coulomb scattering & refers & as term \\
\bottomrule
\end{manualtable}

\Needspace{10\baselineskip}\subsection*{BIOL: Darwinian revolution}
\addcontentsline{toc}{subsection}{BIOL: Darwinian revolution}
\noindent\textit{Predecessor} $T_1$: Natural history of fixed species (Linnaeus, Cuvier, Paley, Owen), c.~1750--1859.\\
\textit{Successor} $T_2$: Darwin's theory of descent with modification by natural selection (Origin of Species, 1859) and its continuation.\par\smallskip
\begin{manualtable}
\textbf{BI01} & \textbf{species}. A kind of organism created as such, fixed, with limits of variation set by its type. & species (a lineage of interbreeding populations) & \textbf{P} partial; $w$\,=\,11; $s$\,=\,1 & H: I16 I20 I22 I23 I30 I31; X: I01 I02 I03 I26; U: I17 \\
\quad BI01a & Felis leo in Linnaeus' Systema Naturae & the lion & refers & as term \\
\quad BI01b & The Galapagos finches as named by Gould (1837) & the finch species & refers & as term \\
\midrule
\textbf{BI02} & \textbf{variety}. A form within a species that departs from the type and does not persist. & variety (an incipient species) & \textbf{P} partial; $w$\,=\,1; $s$\,=\,1 & X: I02 \\
\quad BI02a & Varieties of cultivated plants & varieties & refers & as term \\
\midrule
\textbf{BI03} & \textbf{genus}. A group of species sharing a common character, an order of the natural system. & genus (a group of species descended from a common ancestor) & \textbf{E} exact; $w$\,=\,2; $s$\,=\,1 & H: I04 I31 \\
\quad BI03a & The genus Felis & the genus Felis & refers & as term \\
\midrule
\textbf{BI04} & \textbf{natural system}. The arrangement of organisms in groups subordinate to groups, expressing their true affinities. & genealogical classification & \textbf{P} partial; $w$\,=\,3; $s$\,=\,1 & H: I04 I06; U: I05 \\
\quad BI04a & The natural system sought by de Candolle and Cuvier & genealogical classification & refers & as term \\
\midrule
\textbf{BI05} & \textbf{plan of creation}. The design in the Creator's mind of which the natural system is the expression. & --- & \textbf{T} token-dependent; $w$\,=\,3; $s$\,=\,1/2 & by token \\
\quad BI05a & The plan as the cause of the natural system & \textit{no referent} & fails & U: I05 I11 \\
\quad BI05b & Unity of plan: the shared structure of the members of a great class & nested pattern of shared ancestral structure & refers & H: I06 \\
\midrule
\textbf{BI06} & \textbf{archetype}. Owen (1848): the ideal pattern of which all vertebrate skeletons are modifications. & --- & \textbf{T} token-dependent; $w$\,=\,2; $s$\,=\,1/2 & by token \\
\quad BI06a & The archetype as an idea in the divine mind & \textit{no referent} & fails & U: I07 \\
\quad BI06b & The generalized vertebrate pattern drawn by Owen & the structure of the common ancestor of the vertebrates (Darwin, Origin, ch.~13) & refers & H: I08 \\
\midrule
\textbf{BI07} & \textbf{homology}. The same organ in different animals under every variety of form and function (Owen). & homology (similarity due to common descent) & \textbf{E} exact; $w$\,=\,3; $s$\,=\,1 & H: I08 I09 I10 \\
\quad BI07a & The forelimb bones of man, bat and porpoise & homologous forelimbs & refers & as term \\
\midrule
\textbf{BI08} & \textbf{analogy}. Similarity of function in organs of different structure. & analogy (convergent similarity) & \textbf{E} exact; $w$\,=\,1; $s$\,=\,1 & H: I10 \\
\quad BI08a & Wings of birds and of insects & analogous wings & refers & as term \\
\midrule
\textbf{BI09} & \textbf{adaptation (contrivance)}. The fitting of an organism's parts to its conditions of life, evidence of design (Paley, 1802). & adaptation (product of natural selection) & \textbf{P} partial; $w$\,=\,2; $s$\,=\,1 & H: I15; X: I14 \\
\quad BI09a & Paley's eye & the eye as an adaptation & refers & as term \\
\midrule
\textbf{BI10} & \textbf{centre of creation}. The place where a species was created. & --- & \textbf{T} token-dependent; $w$\,=\,2; $s$\,=\,1/2 & by token \\
\quad BI10a & The place and act of separate creation of a species & \textit{no referent} & fails & U: I17 \\
\quad BI10b & The single region from which a species spread & the area of origin of the species (Darwin's single centres) & refers & H: I16 \\
\midrule
\textbf{BI11} & \textbf{rudimentary organ}. An organ present in reduced form, retained for the completion of the plan or for symmetry. & vestigial organ & \textbf{P} partial; $w$\,=\,2; $s$\,=\,1 & H: I12; U: I11 \\
\quad BI11a & Teeth of foetal whales; mammae of male mammals & vestigial organs & refers & as term \\
\midrule
\textbf{BI12} & \textbf{embryonic resemblance}. The greater resemblance of embryos than of adults within a class (von Baer). & embryonic resemblance & \textbf{E} exact; $w$\,=\,1; $s$\,=\,1 & H: I13 \\
\quad BI12a & Von Baer's unlabelled embryos & embryonic resemblance & refers & as term \\
\midrule
\textbf{BI13} & \textbf{extinction}. The disappearance of species, in Cuvier by revolutions of the globe. & extinction & \textbf{P} partial; $w$\,=\,2; $s$\,=\,1 & H: I30; X: I18 \\
\quad BI13a & Extinction of the mammoth and the Megatherium & extinction & refers & as term \\
\midrule
\textbf{BI14} & \textbf{succession of types}. The resemblance of extinct to living forms in the same region (Clift, Owen). & succession of types & \textbf{E} exact; $w$\,=\,1; $s$\,=\,1 & H: I20 \\
\quad BI14a & Fossil edentates of South America and living armadillos & succession of types & refers & as term \\
\midrule
\textbf{BI15} & \textbf{fossil sequence}. The order of appearance of fossil forms in successive strata. & fossil sequence & \textbf{E} exact; $w$\,=\,1; $s$\,=\,1 & H: I19 \\
\quad BI15a & The faunal succession of William Smith and Cuvier & fossil sequence & refers & as term \\
\midrule
\textbf{BI16} & \textbf{chain of being}. A single linear scale of organisms from lowest to highest, without gaps. & --- & \textbf{N} none; $w$\,=\,1; $s$\,=\,0 & by token \\
\quad BI16a & The scala naturae of Bonnet & \textit{no referent} & fails & U: I21 \\
\midrule
\textbf{BI17} & \textbf{domestic race}. A breed of domestic animal or cultivated plant. & breed & \textbf{E} exact; $w$\,=\,2; $s$\,=\,1 & H: I22 I28 \\
\quad BI17a & Breeds of pigeons & breeds & refers & as term \\
\midrule
\textbf{BI18} & \textbf{hybrid sterility}. The infertility of crosses between species, a provision against their confusion. & hybrid sterility & \textbf{P} partial; $w$\,=\,2; $s$\,=\,1 & H: I23; X: I24 \\
\quad BI18a & Sterility of the mule & hybrid sterility & refers & as term \\
\midrule
\textbf{BI19} & \textbf{struggle for existence}. The checks on increase described by Malthus and de Candolle's war of nature. & struggle for existence & \textbf{E} exact; $w$\,=\,1; $s$\,=\,1 & H: I25 \\
\quad BI19a & Malthus's checks on population & struggle for existence & refers & as term \\
\midrule
\textbf{BI20} & \textbf{balance of nature}. A designed equilibrium in which each species keeps its station (Linnaeus' economy of nature). & the network of ecological relations & \textbf{P} partial; $w$\,=\,1; $s$\,=\,1 & X: I26 \\
\quad BI20a & Linnaeus' Oeconomia naturae & ecological relations & refers & as term \\
\midrule
\textbf{BI21} & \textbf{instinct}. A fixed endowment of each species, as perfect as its structures. & instinct (heritable behaviour subject to selection) & \textbf{P} partial; $w$\,=\,1; $s$\,=\,1 & X: I27 \\
\quad BI21a & The cell-making instinct of the hive bee & instinct & refers & as term \\
\midrule
\textbf{BI22} & \textbf{sport}. A sudden large departure from the parental form. & sport (a large discontinuous variation) & \textbf{E} exact; $w$\,=\,1; $s$\,=\,1 & H: I28 \\
\quad BI22a & Ancon sheep & a large heritable variation & refers & as term \\
\midrule
\textbf{BI23} & \textbf{variation}. Individual differences, confined within the limits of the type. & variation (heritable individual differences) & \textbf{P} partial; $w$\,=\,2; $s$\,=\,1 & H: I22; X: I03 \\
\quad BI23a & Differences among individuals of a litter & variation & refers & as term \\
\midrule
\textbf{BI24} & \textbf{higher and lower organisms}. Position on the scale of perfection. & degree of differentiation and specialization of parts & \textbf{E} exact; $w$\,=\,1; $s$\,=\,1 & H: I29 \\
\quad BI24a & Mammals higher than fishes & degree of differentiation & refers & as term \\
\bottomrule
\end{manualtable}

\Needspace{10\baselineskip}\subsection*{PHYS: Harveian revolution}
\addcontentsline{toc}{subsection}{PHYS: Harveian revolution}
\noindent\textit{Predecessor} $T_1$: Galenic physiology (Galen, On the Natural Faculties, On the Use of Parts; as taught to 1628).\\
\textit{Successor} $T_2$: Harvey's theory of the circulation (De motu cordis, 1628) and its continuation (Malpighi's capillaries, 1661).\par\smallskip
\begin{manualtable}
\textbf{GA01} & \textbf{venous blood}. Nutritive blood made in the liver from chyle and distributed by the veins to be consumed by the parts. & venous blood & \textbf{P} partial; $w$\,=\,5; $s$\,=\,1 & H: I08 I26 I29; X: I01 I02 \\
\quad GA01a & Blood in the vena cava & venous blood & refers & as term \\
\quad GA01b & Blood let from a vein of the arm & venous blood & refers & as term \\
\midrule
\textbf{GA02} & \textbf{arterial blood}. Thin, bright blood mixed with vital spirit, distributed by the arteries. & arterial blood & \textbf{P} partial; $w$\,=\,4; $s$\,=\,1 & H: I08 I23 I29; U: I07 \\
\quad GA02a & Blood from a cut artery & arterial blood & refers & as term \\
\midrule
\textbf{GA03} & \textbf{vital spirit}. Pneuma elaborated in the left ventricle from inspired air and blood; carried by the arteries to vivify the body. & --- & \textbf{N} none; $w$\,=\,4; $s$\,=\,0 & by token \\
\quad GA03a & The vital spirit of the left ventricle & \textit{no referent} & fails & U: I06 I07 I14 I17 \\
\midrule
\textbf{GA04} & \textbf{natural spirit}. Spirit elaborated in the liver and carried by the veins, governing nutrition and growth. & --- & \textbf{N} none; $w$\,=\,1; $s$\,=\,0 & by token \\
\quad GA04a & The natural spirit of the liver & \textit{no referent} & fails & U: I20 \\
\midrule
\textbf{GA05} & \textbf{animal spirit}. Spirit elaborated in the brain from vital spirit, carried by the nerves, producing sensation and motion. & --- & \textbf{T} token-dependent; $w$\,=\,2; $s$\,=\,1/2 & by token \\
\quad GA05a & Spirit elaborated in the rete mirabile and ventricles of the brain & \textit{no referent} & fails & U: I17 \\
\quad GA05b & What passes along the nerves from the brain to move the muscles & nerve conduction & refers & H: I18 \\
\midrule
\textbf{GA06} & \textbf{liver}. The origin of the veins and the organ that makes the blood. & liver & \textbf{P} partial; $w$\,=\,3; $s$\,=\,1 & X: I01 I22; U: I20 \\
\quad GA06a & The liver dissected & liver & refers & as term \\
\midrule
\textbf{GA07} & \textbf{heart}. The origin of the arteries and seat of the innate heat; attracts blood in diastole. & heart (a muscular pump) & \textbf{P} partial; $w$\,=\,4; $s$\,=\,1 & H: I30; U: I03 I06 I13 \\
\quad GA07a & The heart dissected and observed beating & heart & refers & as term \\
\midrule
\textbf{GA08} & \textbf{septal pores}. Invisible pores in the septum through which blood passes from the right to the left ventricle. & --- & \textbf{N} none; $w$\,=\,1; $s$\,=\,0 & by token \\
\quad GA08a & The pores of the interventricular septum & \textit{no referent} & fails & U: I03 \\
\midrule
\textbf{GA09} & \textbf{arterial vein}. The vessel carrying blood from the right ventricle to the lungs to nourish them. & pulmonary artery & \textbf{P} partial; $w$\,=\,2; $s$\,=\,1 & H: I04; X: I27 \\
\quad GA09a & Vena arteriosa & pulmonary artery & refers & as term \\
\midrule
\textbf{GA10} & \textbf{venous artery}. The vessel carrying air from the lungs to the left ventricle and sooty vapours back. & pulmonary vein & \textbf{P} partial; $w$\,=\,2; $s$\,=\,1 & X: I05; U: I15 \\
\quad GA10a & Arteria venosa & pulmonary veins & refers & as term \\
\midrule
\textbf{GA11} & \textbf{diastole}. The dilation of the heart, its active motion, by which it attracts blood. & diastole (relaxation and filling) & \textbf{P} partial; $w$\,=\,1; $s$\,=\,1 & U: I09 \\
\quad GA11a & Diastole observed in a vivisected heart & diastole & refers & as term \\
\midrule
\textbf{GA12} & \textbf{systole}. The contraction of the heart. & systole (active contraction and ejection) & \textbf{E} exact; $w$\,=\,2; $s$\,=\,1 & H: I10 I30 \\
\quad GA12a & Systole observed in a vivisected heart & systole & refers & as term \\
\midrule
\textbf{GA13} & \textbf{pulse}. The dilation and contraction of the arteries, produced by a pulsific faculty conveyed from the heart through their coats. & arterial pulse & \textbf{P} partial; $w$\,=\,2; $s$\,=\,1 & H: I24; U: I11 \\
\quad GA13a & The pulse felt at the wrist & arterial pulse & refers & as term \\
\midrule
\textbf{GA14} & \textbf{innate heat}. The heat seated in the heart, source of life, moderated by respiration. & --- & \textbf{T} token-dependent; $w$\,=\,2; $s$\,=\,1/2 & by token \\
\quad GA14a & Innate heat as the principle of life seated in the heart & \textit{no referent} & fails & U: I13 \\
\quad GA14b & The warmth of the living body, felt to the touch & body heat (produced by metabolism) & refers & H: I28 \\
\midrule
\textbf{GA15} & \textbf{sooty vapours}. Waste vapours produced by the innate heat in the heart and expelled through the lungs. & --- & \textbf{T} token-dependent; $w$\,=\,2; $s$\,=\,1/2 & by token \\
\quad GA15a & Vapours passing back from the left ventricle through the venous artery & \textit{no referent} & fails & U: I15 \\
\quad GA15b & Waste carried away by expired air & carbon dioxide and water vapour in expired air & refers & H: I16 \\
\midrule
\textbf{GA16} & \textbf{chyle}. The product of digestion carried from the intestines to the liver. & chyle & \textbf{P} partial; $w$\,=\,2; $s$\,=\,1 & X: I01 I22 \\
\quad GA16a & The milky juice of digestion & chyle & refers & as term \\
\midrule
\textbf{GA17} & \textbf{valves of the veins}. Membranes in the veins that retard the flow of blood to the extremities (Fabricius, 1603). & venous valves & \textbf{P} partial; $w$\,=\,1; $s$\,=\,1 & X: I12 \\
\quad GA17a & Fabricius' ostiola & venous valves & refers & as term \\
\midrule
\textbf{GA18} & \textbf{faculties}. The attractive, retentive, alterative and expulsive powers of the parts. & --- & \textbf{N} none; $w$\,=\,3; $s$\,=\,0 & by token \\
\quad GA18a & The attractive faculty of the heart and of the parts & \textit{no referent} & fails & U: I09 I11 I25 \\
\midrule
\textbf{GA19} & \textbf{nutrition}. The assimilation by each part of the blood it attracts and consumes. & nutrition & \textbf{P} partial; $w$\,=\,4; $s$\,=\,1 & H: I26; X: I02 I27; U: I25 \\
\quad GA19a & Nourishment of the flesh & nutrition & refers & as term \\
\midrule
\textbf{GA20} & \textbf{lungs}. Organs that cool the heart and supply it with air. & lungs & \textbf{P} partial; $w$\,=\,4; $s$\,=\,1 & H: I04; X: I05 I27; U: I14 \\
\quad GA20a & The lungs dissected & lungs & refers & as term \\
\midrule
\textbf{GA21} & \textbf{respiration}. The drawing in of air to cool the innate heat and to supply the matter of vital spirit. & respiration & \textbf{P} partial; $w$\,=\,2; $s$\,=\,1 & H: I16; U: I14 \\
\quad GA21a & Breathing & respiration & refers & as term \\
\midrule
\textbf{GA22} & \textbf{anastomoses}. Connections between arteries and veins at their extremities, through which some blood and spirit pass. & capillary connections & \textbf{E} exact; $w$\,=\,1; $s$\,=\,1 & H: I21 \\
\quad GA22a & Galen's anastomoses & capillary beds & refers & as term \\
\midrule
\textbf{GA23} & \textbf{portal vein}. The vein carrying chyle from the intestines to the liver. & hepatic portal vein & \textbf{P} partial; $w$\,=\,1; $s$\,=\,1 & X: I22 \\
\quad GA23a & Vena portae & hepatic portal vein & refers & as term \\
\midrule
\textbf{GA24} & \textbf{vena cava}. The great vein carrying blood from the liver to the body. & vena cava & \textbf{P} partial; $w$\,=\,1; $s$\,=\,1 & X: I02 \\
\quad GA24a & Vena cava & vena cava & refers & as term \\
\midrule
\textbf{GA25} & \textbf{aorta}. The great artery from the left ventricle. & aorta & \textbf{E} exact; $w$\,=\,2; $s$\,=\,1 & H: I10 I23 \\
\quad GA25a & Aorta & aorta & refers & as term \\
\midrule
\textbf{GA26} & \textbf{rete mirabile}. A network of arteries at the base of the brain in which vital spirit is elaborated into animal spirit. & --- & \textbf{T} token-dependent; $w$\,=\,2; $s$\,=\,1/2 & by token \\
\quad GA26a & The rete mirabile of the human brain & \textit{no referent} & fails & U: I17 \\
\quad GA26b & The rete mirabile dissected in the ox & the carotid rete of ruminants & refers & H: I19 \\
\bottomrule
\end{manualtable}

\Needspace{10\baselineskip}\subsection*{OPTI: Optical revolution}
\addcontentsline{toc}{subsection}{OPTI: Optical revolution}
\noindent\textit{Predecessor} $T_1$: Newtonian emission optics (Opticks, 1704) and its development (Laplace, Biot, Malus), 1704--c.~1815.\\
\textit{Successor} $T_2$: The wave theory of light (Young 1801--1803; Fresnel 1815--1823).\par\smallskip
\begin{manualtable}
\textbf{OP01} & \textbf{ray}. The least part of light, a stream of corpuscles moving in a straight line in a uniform medium. & ray (normal to the wave fronts) & \textbf{P} partial; $w$\,=\,5; $s$\,=\,1 & H: I05 I21 I22; R: I02; X: I12 \\
\quad OP01a & A sunbeam admitted through a hole in the shutter & a narrow beam of light & refers & as term \\
\quad OP01b & Rays of different colours separated by a prism & rays of different wavelengths & refers & as term \\
\midrule
\textbf{OP02} & \textbf{light corpuscle}. A small body emitted by luminous bodies and acted on by forces of matter. & --- & \textbf{N} none; $w$\,=\,6; $s$\,=\,0 & by token \\
\quad OP02a & The particles of Query 29 of the Opticks & \textit{no referent} & fails & U: I01 I03 I07 I11 I17 I20 \\
\midrule
\textbf{OP03} & \textbf{fits of easy reflection and transmission}. Alternating dispositions of a ray, recurring at equal intervals, to be reflected or transmitted at a surface. & --- & \textbf{T} token-dependent; $w$\,=\,4; $s$\,=\,1/2 & by token \\
\quad OP03a & Fits as dispositions of the corpuscles themselves & \textit{no referent} & fails & U: I07 \\
\quad OP03b & The interval of the fits measured from Newton's rings, from which Young (1802) computed wavelengths & half the wavelength of the light in the film & refers & H: I08 I09 I25 \\
\midrule
\textbf{OP04} & \textbf{colour (of a ray)}. An original and unchangeable property of a ray, connected with its degree of refrangibility. & wavelength (frequency) of the light & \textbf{E} exact; $w$\,=\,4; $s$\,=\,1 & H: I06 I09 I23 I24 \\
\quad OP04a & Red and violet light of the prismatic spectrum & long and short wavelengths & refers & as term \\
\midrule
\textbf{OP05} & \textbf{refrangibility}. The disposition of a ray to be refracted more or less, fixed for each colour. & dependence of refractive index on wavelength & \textbf{P} partial; $w$\,=\,4; $s$\,=\,1 & H: I05 I06 I22; U: I03 \\
\quad OP05a & Greater refrangibility of violet than red & larger index at shorter wavelengths & refers & as term \\
\midrule
\textbf{OP06} & \textbf{velocity of light in a medium}. Greater in a denser medium, because the corpuscles are attracted by it. & phase velocity of light in the medium & \textbf{P} partial; $w$\,=\,1; $s$\,=\,1 & X: I04 \\
\quad OP06a & Velocity of light in water & speed of light in water & refers & as term \\
\midrule
\textbf{OP07} & \textbf{inflection}. The bending of rays near the edges of bodies by forces of the edge. & diffraction & \textbf{P} partial; $w$\,=\,2; $s$\,=\,1 & H: I10; U: I11 \\
\quad OP07a & Grimaldi's and Newton's fringes at the edge of a shadow & diffraction fringes & refers & as term \\
\midrule
\textbf{OP08} & \textbf{sides of rays}. Original differences between the sides of a ray, which account for the unusual refraction of Iceland crystal. & --- & \textbf{T} token-dependent; $w$\,=\,2; $s$\,=\,1/2 & by token \\
\quad OP08a & Sides as properties of the corpuscles & \textit{no referent} & fails & U: I14 \\
\quad OP08b & The property of a beam emerging from Iceland crystal that fixes how it is refracted by a second crystal & polarization (the orientation of the transverse vibration) & refers & H: I13 \\
\midrule
\textbf{OP09} & \textbf{polarization}. Malus (1808--1810): the arrangement of the axes of the light particles by reflection or double refraction. & polarization & \textbf{P} partial; $w$\,=\,3; $s$\,=\,1 & H: I15 I16; U: I17 \\
\quad OP09a & Light reflected from a window and viewed through Iceland crystal (Malus) & polarized light & refers & as term \\
\midrule
\textbf{OP10} & \textbf{double refraction}. The division of a beam into two by Iceland crystal. & birefringence & \textbf{E} exact; $w$\,=\,2; $s$\,=\,1 & H: I13 I16 \\
\quad OP10a & Iceland crystal & birefringence of calcite & refers & as term \\
\midrule
\textbf{OP11} & \textbf{Newton's rings}. Coloured rings formed in a thin film of air between two glasses. & thin-film interference rings & \textbf{E} exact; $w$\,=\,3; $s$\,=\,1 & H: I08 I24 I25 \\
\quad OP11a & Rings between a lens and a flat glass & interference rings & refers & as term \\
\midrule
\textbf{OP12} & \textbf{aberration of starlight}. The annual displacement of stars due to the Earth's motion and the finite speed of light (Bradley, 1729). & aberration & \textbf{E} exact; $w$\,=\,1; $s$\,=\,1 & H: I18 \\
\quad OP12a & Bradley's observations of gamma Draconis & aberration & refers & as term \\
\midrule
\textbf{OP13} & \textbf{white light}. A heterogeneous mixture of rays of all colours. & light of many wavelengths superposed & \textbf{E} exact; $w$\,=\,2; $s$\,=\,1 & H: I05 I23 \\
\quad OP13a & Sunlight & sunlight & refers & as term \\
\midrule
\textbf{OP14} & \textbf{homogeneal light}. Light of one colour and one refrangibility. & monochromatic light & \textbf{E} exact; $w$\,=\,2; $s$\,=\,1 & H: I06 I22 \\
\quad OP14a & Light from one part of the prismatic spectrum & nearly monochromatic light & refers & as term \\
\midrule
\textbf{OP15} & \textbf{reflection}. The return of rays from a surface by a repulsive power acting at a small distance from it. & reflection & \textbf{P} partial; $w$\,=\,3; $s$\,=\,1 & H: I15 I21; U: I20 \\
\quad OP15a & Reflection from glass & reflection & refers & as term \\
\midrule
\textbf{OP16} & \textbf{geometrical shadow}. The region a straight-line propagation leaves dark behind an obstacle. & geometrical shadow & \textbf{P} partial; $w$\,=\,2; $s$\,=\,1 & H: I10; X: I12 \\
\quad OP16a & The shadow of a small disk & geometrical shadow & refers & as term \\
\midrule
\textbf{OP17} & \textbf{intensity}. The density of rays. & intensity & \textbf{E} exact; $w$\,=\,1; $s$\,=\,1 & H: I19 \\
\quad OP17a & Brightness of a candle at different distances & intensity & refers & as term \\
\midrule
\textbf{OP18} & \textbf{prismatic spectrum}. The oblong image formed by a prism from sunlight. & dispersion spectrum & \textbf{E} exact; $w$\,=\,1; $s$\,=\,1 & H: I23 \\
\quad OP18a & Newton's spectrum in the experimentum crucis & dispersion spectrum & refers & as term \\
\bottomrule
\end{manualtable}

\Needspace{10\baselineskip}\subsection*{MIAS: Bacteriological revolution}
\addcontentsline{toc}{subsection}{MIAS: Bacteriological revolution}
\noindent\textit{Predecessor} $T_1$: Miasmatic and zymotic theories of epidemic disease (Chadwick, Southwood Smith, Farr; Liebig's chemical theory of ferments), c.~1840--1875.\\
\textit{Successor} $T_2$: The germ theory of disease and fermentation (Pasteur, Lister, Koch), 1857--1890, and its continuation.\par\smallskip
\begin{manualtable}
\textbf{MI01} & \textbf{miasma}. A noxious effluvium from decaying organic matter, marshes and filth, which produces epidemic disease when breathed. & --- & \textbf{T} token-dependent; $w$\,=\,5; $s$\,=\,2/3 & by token \\
\quad MI01a & The poisonous emanation of putrefying matter as such & \textit{no referent} & fails & U: I01 I14 I19 \\
\quad MI01b & The cause of intermittent fevers in marshy districts & infection by Plasmodium, carried by mosquitoes breeding in marshes (Laveran, 1880; Ross, 1897) & refers & H: I02 \\
\quad MI01c & The cause of cholera in low-lying districts with foul air and water (Farr) & Vibrio cholerae carried in contaminated water & refers & R: I31 \\
\midrule
\textbf{MI02} & \textbf{zymotic disease}. Farr's class of epidemic, endemic and contagious diseases, produced by a process in the blood analogous to fermentation. & infectious disease & \textbf{P} partial; $w$\,=\,2; $s$\,=\,1 & H: I24; U: I03 \\
\quad MI02a & Smallpox, measles, scarlatina, typhus and cholera in Farr's nosology & these infectious diseases & refers & as term \\
\midrule
\textbf{MI03} & \textbf{ferment}. An agent that induces fermentation or putrefaction in another substance. & --- & \textbf{T} token-dependent; $w$\,=\,4; $s$\,=\,2/3 & by token \\
\quad MI03a & Yeast in alcoholic fermentation (Cagniard-Latour, Schwann, 1837; Pasteur, 1857--1860) & yeast cells & refers & H: I07 \\
\quad MI03b & Diastase, extracted from malt (Payen and Persoz, 1833) & amylase (an enzyme) & refers & H: I29 \\
\quad MI03c & Liebig: a body in decomposition communicating its molecular motion to the fermentable substance & \textit{no referent} & fails & U: I03 I08 \\
\midrule
\textbf{MI04} & \textbf{contagion}. Communication of a disease from the sick to the healthy by contact or by objects they have touched. & transmission of infection & \textbf{E} exact; $w$\,=\,2; $s$\,=\,1 & H: I04 I27 \\
\quad MI04a & Communication of smallpox to attendants and by clothing & transmission of smallpox & refers & as term \\
\midrule
\textbf{MI05} & \textbf{contagium}. The specific matter by which a contagious disease is communicated. & --- & \textbf{T} token-dependent; $w$\,=\,4; $s$\,=\,2/3 & by token \\
\quad MI05a & Variolous matter used in inoculation & variola virus & refers & H: I04 I24 \\
\quad MI05b & Henle's contagium animatum (1840): a living agent & pathogenic microorganisms & refers & H: I05 \\
\quad MI05c & A chemical poison that multiplies by inducing decomposition in the blood & \textit{no referent} & fails & U: I06 \\
\midrule
\textbf{MI06} & \textbf{putrefaction}. Spontaneous chemical decomposition of dead organic matter, initiated by oxygen. & putrefaction (decomposition by microorganisms) & \textbf{P} partial; $w$\,=\,2; $s$\,=\,1 & H: I30; X: I09 \\
\quad MI06a & Putrefaction of meat broth & microbial decomposition & refers & as term \\
\midrule
\textbf{MI07} & \textbf{spontaneous generation}. Origin of microscopic organisms in infusions without parent organisms. & --- & \textbf{T} token-dependent; $w$\,=\,2; $s$\,=\,1/2 & by token \\
\quad MI07a & The doctrine that living beings arise from non-living matter & \textit{no referent} & fails & U: I10 \\
\quad MI07b & The organisms appearing in boiled hay infusions (Pouchet, 1859) & organisms grown from heat-resistant spores (Cohn; Tyndall, 1876--1877) & refers & H: I11 \\
\midrule
\textbf{MI08} & \textbf{epidemic constitution}. A state of the atmosphere favouring a particular epidemic (Sydenham). & --- & \textbf{N} none; $w$\,=\,2; $s$\,=\,0 & by token \\
\quad MI08a & The atmospheric constitution of an epidemic year & \textit{no referent} & fails & U: I12 I21 \\
\midrule
\textbf{MI09} & \textbf{intermittent fever}. A fever recurring at regular intervals, endemic in marshy districts. & malaria & \textbf{E} exact; $w$\,=\,2; $s$\,=\,1 & H: I02 I28 \\
\quad MI09a & Tertian and quartan fevers & malaria (tertian and quartan forms) & refers & as term \\
\midrule
\textbf{MI10} & \textbf{continued fever}. A fever without regular intermission, treated as one disease varying in form. & typhus $\mid$ typhoid fever (pd) & \textbf{P} partial; $w$\,=\,1; $s$\,=\,1 & X: I25 \\
\quad MI10a & The continued fever of the London Fever Hospital before 1850 & typhus or typhoid (partial denotation) & refers & as term \\
\midrule
\textbf{MI11} & \textbf{filth}. Accumulated organic refuse, excrement and decaying matter, the source of miasma. & organic refuse (a source of pathogens and of their vectors) & \textbf{P} partial; $w$\,=\,2; $s$\,=\,1 & H: I13; U: I14 \\
\quad MI11a & Cesspools and refuse of the courts described in Chadwick's Report (1842) & sources of faecal contamination & refers & as term \\
\midrule
\textbf{MI12} & \textbf{sewer gas}. Gas escaping from drains, a cause of typhoid and other fevers. & --- & \textbf{N} none; $w$\,=\,1; $s$\,=\,0 & by token \\
\quad MI12a & Sewer gas as the cause of typhoid & \textit{no referent} & fails & U: I15 \\
\midrule
\textbf{MI13} & \textbf{cholera poison}. The specific poison producing cholera. & --- & \textbf{T} token-dependent; $w$\,=\,2; $s$\,=\,1/2 & by token \\
\quad MI13a & Snow (1849, 1855): morbid matter that multiplies in the gut and passes in the evacuations into water & Vibrio cholerae & refers & H: I16 \\
\quad MI13b & Pettenkofer: a poison that must mature in porous soil before it can cause disease & \textit{no referent} & fails & U: I17 \\
\midrule
\textbf{MI14} & \textbf{disinfectant}. An agent that destroys miasma and removes offensive smells. & disinfectant (an agent killing microorganisms) & \textbf{P} partial; $w$\,=\,2; $s$\,=\,1 & H: I20; U: I19 \\
\quad MI14a & Chloride of lime & hypochlorite disinfectant & refers & as term \\
\midrule
\textbf{MI15} & \textbf{vaccination}. Inoculation with cowpox matter to protect against smallpox (Jenner, 1798). & vaccination & \textbf{E} exact; $w$\,=\,1; $s$\,=\,1 & H: I22 \\
\quad MI15a & Jennerian vaccination & vaccination with cowpox virus & refers & as term \\
\midrule
\textbf{MI16} & \textbf{hospital gangrene}. A spreading destruction of wounds in crowded hospital wards. & wound infection (necrotizing) & \textbf{E} exact; $w$\,=\,1; $s$\,=\,1 & H: I23 \\
\quad MI16a & Hospital gangrene in military hospitals & necrotizing wound infection & refers & as term \\
\midrule
\textbf{MI17} & \textbf{puerperal fever}. Fever following childbirth, often epidemic in lying-in wards. & puerperal sepsis & \textbf{P} partial; $w$\,=\,2; $s$\,=\,1 & H: I20; U: I21 \\
\quad MI17a & Puerperal fever in the Vienna maternity clinic (Semmelweis, 1847) & puerperal sepsis & refers & as term \\
\midrule
\textbf{MI18} & \textbf{infection}. The taking of a disease from a source, by air, contact or objects. & infection (entry and multiplication of a pathogen) & \textbf{E} exact; $w$\,=\,2; $s$\,=\,1 & H: I27 I30 \\
\quad MI18a & Infection by fomites & infection & refers & as term \\
\midrule
\textbf{MI19} & \textbf{predisposing cause}. A condition of the body or of life that renders a person liable to an epidemic. & host susceptibility & \textbf{E} exact; $w$\,=\,1; $s$\,=\,1 & H: I26 \\
\quad MI19a & Poverty, intemperance, debility & factors of susceptibility & refers & as term \\
\midrule
\textbf{MI20} & \textbf{fermentation}. Decomposition of a substance induced by a ferment. & fermentation (microbial or enzymatic metabolism) & \textbf{P} partial; $w$\,=\,2; $s$\,=\,1 & H: I07; U: I08 \\
\quad MI20a & Alcoholic fermentation of grape juice & alcoholic fermentation by yeast & refers & as term \\
\midrule
\textbf{MI21} & \textbf{cholera}. Asiatic cholera, an epidemic disease with violent purging. & cholera (infection by Vibrio cholerae) & \textbf{P} partial; $w$\,=\,3; $s$\,=\,1 & H: I16 I18; R: I31 \\
\quad MI21a & The London epidemics of 1849 and 1854 & cholera & refers & as term \\
\bottomrule
\end{manualtable}

\Needspace{10\baselineskip}\subsection*{HERE: Mendelian revolution}
\addcontentsline{toc}{subsection}{HERE: Mendelian revolution}
\noindent\textit{Predecessor} $T_1$: Pre-Mendelian theories of heredity: blending inheritance, Darwin's pangenesis (1868), Galton's ancestral heredity, the inheritance of acquired characters, c.~1850--1900.\\
\textit{Successor} $T_2$: Mendelian genetics (Mendel 1866, rediscovered 1900) and the chromosome theory (Morgan, 1910--1915).\par\smallskip
\begin{manualtable}
\textbf{HR01} & \textbf{gemmule}. Darwin (1868): minute particles thrown off by every cell, circulating through the body and collecting in the reproductive elements; may lie dormant for generations. & --- & \textbf{T} token-dependent; $w$\,=\,2; $s$\,=\,1/2 & by token \\
\quad HR01a & Gemmules as particles shed by the cells of the body & \textit{no referent} & fails & U: I01 \\
\quad HR01b & Dormant gemmules transmitted unexpressed through generations, to which Darwin ascribed reversion & recessive alleles carried unexpressed & refers & R: I06 \\
\midrule
\textbf{HR02} & \textbf{blending inheritance}. The mixing of the parents' hereditary natures in the offspring, so that their characters are averaged. & intermediate inheritance of quantitative characters & \textbf{P} partial; $w$\,=\,5; $s$\,=\,1 & R: I22; X: I02 I03 I11 I23 \\
\quad HR02a & Children of tall and short parents intermediate in stature & quantitative (polygenic) inheritance & refers & as term \\
\quad HR02b & Jenkin (1867): the dilution of a rare variation by crossing & the frequency of a rare variant in a population & refers & as term \\
\midrule
\textbf{HR03} & \textbf{prepotency}. The stronger power of one parent, or one breed, to transmit its characters to the offspring. & dominance & \textbf{E} exact; $w$\,=\,1; $s$\,=\,1 & H: I04 \\
\quad HR03a & Darwin on prepotent breeds in crosses & dominance of an allele & refers & as term \\
\midrule
\textbf{HR04} & \textbf{reversion}. The reappearance of an ancestral character after one or more generations. & reappearance of a recessive or recombined phenotype & \textbf{P} partial; $w$\,=\,2; $s$\,=\,1 & H: I05; R: I06 \\
\quad HR04a & Crosses of distinct breeds of pigeon yielding the blue colouring of the rock pigeon & recombination of complementary alleles & refers & as term \\
\midrule
\textbf{HR05} & \textbf{use-inheritance}. The transmission to offspring of characters acquired by use or disuse. & --- & \textbf{N} none; $w$\,=\,1; $s$\,=\,0 & by token \\
\quad HR05a & Inheritance of effects of use and disuse & \textit{no referent} & fails & U: I07 \\
\midrule
\textbf{HR06} & \textbf{telegony}. The influence of a previous sire on offspring of a female by later sires. & --- & \textbf{N} none; $w$\,=\,1; $s$\,=\,0 & by token \\
\quad HR06a & Lord Morton's mare and quagga (1820) & \textit{no referent} & fails & U: I08 \\
\midrule
\textbf{HR07} & \textbf{pure breed}. A breed or race that reproduces its characters without variation. & line homozygous for the relevant characters & \textbf{P} partial; $w$\,=\,3; $s$\,=\,1 & H: I05 I09; X: I14 \\
\quad HR07a & Mendel's true-breeding pea varieties & homozygous lines & refers & as term \\
\midrule
\textbf{HR08} & \textbf{hybrid}. The offspring of a cross between distinct breeds or species. & heterozygote (first-generation hybrid) & \textbf{P} partial; $w$\,=\,2; $s$\,=\,1 & H: I10; X: I11 \\
\quad HR08a & Hybrid peas of the first generation & heterozygotes & refers & as term \\
\midrule
\textbf{HR09} & \textbf{sport}. A sudden marked variation appearing in a single individual. & mutation (with visible effect) & \textbf{E} exact; $w$\,=\,1; $s$\,=\,1 & H: I12 \\
\quad HR09a & Sports in cultivated plants & mutations & refers & as term \\
\midrule
\textbf{HR10} & \textbf{individual variation}. The differences among individuals of one breed or species. & heritable (genetic) variation $\mid$ non-heritable variation induced by the environment (pd) & \textbf{P} partial; $w$\,=\,4; $s$\,=\,1 & R: I13; X: I03 I14 I23 \\
\quad HR10a & Differences in seed size within a bean variety & genetic or environmental variation (partial denotation; Johannsen, 1909) & refers & as term \\
\midrule
\textbf{HR11} & \textbf{character}. A feature of an organism transmitted in heredity. & phenotypic trait & \textbf{P} partial; $w$\,=\,6; $s$\,=\,1 & H: I04 I10 I20; R: I22; X: I02; U: I07 \\
\quad HR11a & Seed colour and shape in peas & phenotypic traits & refers & as term \\
\midrule
\textbf{HR12} & \textbf{hereditary tendency}. The disposition, received from the parents, to develop certain characters. & genotype & \textbf{E} exact; $w$\,=\,1; $s$\,=\,1 & H: I15 \\
\quad HR12a & The tendencies of an individual to produce a character in its offspring & genotype & refers & as term \\
\midrule
\textbf{HR13} & \textbf{ancestral heredity}. Galton (1897): each ancestor contributes a fixed fraction of the heritage, one half from the parents, one quarter from the grandparents, and so on. & --- & \textbf{T} token-dependent; $w$\,=\,2; $s$\,=\,1/2 & by token \\
\quad HR13a & The law as a statement of contributions from ancestral germs & \textit{no referent} & fails & U: I17 \\
\quad HR13b & The measured regression of offspring on parents in stature & regression predicted by the additive polygenic model (Fisher, 1918) & refers & H: I18 \\
\midrule
\textbf{HR14} & \textbf{regression}. The tendency of offspring of exceptional parents to be less exceptional (Galton, 1886). & regression to the mean & \textbf{E} exact; $w$\,=\,2; $s$\,=\,1 & H: I16 I18 \\
\quad HR14a & Galton's stature data & regression to the mean & refers & as term \\
\midrule
\textbf{HR15} & \textbf{germ-plasm}. Weismann: the hereditary substance, continuous across generations and not affected by the body. & germ line (chromosomal hereditary material) & \textbf{E} exact; $w$\,=\,1; $s$\,=\,1 & H: I19 \\
\quad HR15a & Weismann's germ-plasm (1892) & germ-line hereditary material & refers & as term \\
\midrule
\textbf{HR16} & \textbf{latent character}. A character carried but not expressed in an individual, and expressed in its descendants. & recessive or sex-limited character & \textbf{E} exact; $w$\,=\,1; $s$\,=\,1 & H: I20 \\
\quad HR16a & Characters latent in one sex and expressed in the other & sex-limited characters & refers & as term \\
\quad HR16b & Characters skipping a generation & recessive characters & refers & as term \\
\midrule
\textbf{HR17} & \textbf{reproductive element}. The ovule or pollen grain, egg or spermatozoon, by which each parent contributes to the offspring. & gamete & \textbf{P} partial; $w$\,=\,3; $s$\,=\,1 & H: I15 I21; U: I01 \\
\quad HR17a & Pollen and ovules of peas & gametes & refers & as term \\
\bottomrule
\end{manualtable}

